\subsection{Overall Performance Comparison}

Table~\ref{tab:overall_comparison} presents aggregate results across all 50 states for both methods.

\begin{table}[h]
\centering
\begin{tabular}{lcccc}
\toprule
\textbf{Method} & \textbf{Configs} & \textbf{Success Rate} & \textbf{Avg Runtime} & \textbf{Total Time} \\
\midrule
N-way          & 880 & 47.5\% (418/880) & 3.68s & 0.90h \\
Recursive      & 880 & \textbf{48.3\%} (425/880) & 11.33s & 2.77h \\
\midrule
Difference     & --  & +0.8 pts & +7.65s (+208\%) & +1.87h \\
\bottomrule
\end{tabular}
\caption{Overall performance comparison across 44 multi-district states, 20 configurations per method per state (5 weight factors $\times$ 4 thresholds). Success rate = fraction achieving proportional MM district target. Single-district states (Alaska, Delaware, Montana, North Dakota, South Dakota, Vermont, Wyoming) excluded from analysis.}
\label{tab:overall_comparison}
\end{table}

\textbf{Finding 1: Nearly Identical Success Rates} — Recursive bisection achieves 48.3\% success rate vs n-way's 47.5\%, a difference of only 0.8 percentage points. This small difference contradicts expectations that architectural choices substantially impact VRA performance.

\textbf{Finding 2: Dramatic Speed Difference} — N-way executes in 3.68s average vs recursive's 11.33s, making n-way \textbf{67.5\% faster} (208\% slower for recursive). Total runtime for complete ablation: 0.90 hours (n-way) vs 2.77 hours (recursive). Speed advantage scales linearly with number of redistricting runs.

\subsection{Statistical Significance Testing}

To rigorously assess whether the 0.8 percentage point difference represents meaningful performance gap or sampling noise, we conduct paired statistical tests on state-level success rates.

\textbf{Paired t-test}: For each of 44 states, we compute average success rate across 20 configurations for each method, yielding paired samples. Test statistic: $t = -0.480$, $p = 0.634$. At conventional $\alpha = 0.05$ significance level, we \textbf{fail to reject the null hypothesis of equal means}. The observed difference is not statistically significant.

\textbf{Effect size (Cohen's d)}: We compute Cohen's d = -0.018, indicating \textbf{negligible effect size} by conventional standards (small: 0.2, medium: 0.5, large: 0.8). The negative sign indicates recursive's slight advantage, but magnitude is trivial.

\textbf{Interpretation}: The 0.8 percentage point difference (425 vs 418 successes out of 880 trials) likely reflects random variation rather than systematic performance gap. With negligible effect size and non-significant p-value, we conclude \textbf{methods perform equivalently for VRA compliance}.

\begin{figure}[h]
\centering
\includegraphics[width=0.7\textwidth]{figures/fig1_overall_success_rates.png}
\caption{Overall VRA compliance success rates for n-way and recursive bisection methods across 44 multi-district states with 20 parameter configurations each (880 runs per method). The difference of 0.8 percentage points is not statistically significant (paired t-test, $p=0.634$, Cohen's $d=-0.018$).}
\label{fig:overall_success}
\end{figure}

\subsection{Parameter Sensitivity Analysis}

While overall performance is equivalent, methods may respond differently to edge-weighting parameters.

\subsubsection{Weight Factor Sensitivity}

Table~\ref{tab:weight_sensitivity} shows success rates by weight factor for both methods.

\begin{table}[h]
\centering
\begin{tabular}{ccccc}
\toprule
\textbf{Weight ($\alpha$)} & \textbf{N-way} & \textbf{Recursive} & \textbf{Difference} & \textbf{Winner} \\
\midrule
1$\times$  (baseline) & 43.2\% (76/176) & 43.2\% (76/176) & 0.0 pts & Tie \\
5$\times$            & 46.6\% (82/176) & 48.3\% (85/176) & +1.7 pts & Recursive \\
10$\times$           & 48.3\% (85/176) & 47.7\% (84/176) & -0.6 pts & N-way \\
50$\times$           & 50.0\% (88/176) & 50.6\% (89/176) & +0.6 pts & Recursive \\
100$\times$          & 49.4\% (87/176) & \textbf{51.7\%} (91/176) & +2.3 pts & Recursive \\
\bottomrule
\end{tabular}
\caption{Success rate by weight factor (collapsed across all thresholds). Each weight factor tested on 44 states $\times$ 4 thresholds = 176 configurations. Winner indicates which method achieves higher success rate for that weight factor.}
\label{tab:weight_sensitivity}
\end{table}

\textbf{Baseline equivalence}: At $\alpha=1$ (no edge-weighting), methods achieve identical 43.2\% success, confirming that performance differences emerge from edge-weight interaction, not base algorithm.

\textbf{Divergence at high weights}: Gap widens at extreme weights. At $\alpha=100$, recursive leads by 2.3 points (51.7\% vs 49.4\%). This suggests recursive's hierarchical structure better leverages strong edge weights—each binary split fully respects edge weights, while n-way's simultaneous $k$-way optimization may dilute edge weight influence.

\textbf{Monotonic improvement}: Both methods show generally increasing success rates from $\alpha=1$ to $\alpha=100$ (with minor non-monotonicity). This confirms edge-weighting improves VRA compliance for both architectures.

\subsubsection{Minority Threshold Sensitivity}

Table~\ref{tab:threshold_sensitivity} shows success rates by minority threshold.

\begin{table}[h]
\centering
\begin{tabular}{ccccc}
\toprule
\textbf{Threshold ($\tau$)} & \textbf{N-way} & \textbf{Recursive} & \textbf{Difference} & \textbf{Winner} \\
\midrule
0.40 & \textbf{49.1\%} (108/220) & 47.3\% (104/220) & -1.8 pts & N-way \\
0.45 & 48.6\% (107/220) & 48.6\% (107/220) & 0.0 pts & Tie \\
0.50 & 46.4\% (102/220) & \textbf{50.9\%} (112/220) & +4.5 pts & Recursive \\
0.55 & 45.9\% (101/220) & 46.4\% (102/220) & +0.5 pts & Recursive \\
\bottomrule
\end{tabular}
\caption{Success rate by minority threshold (collapsed across all weight factors). Each threshold tested on 44 states $\times$ 5 weights = 220 configurations.}
\label{tab:threshold_sensitivity}
\end{table}

\textbf{Opposite threshold preferences}: N-way optimizes at low threshold ($\tau=0.40$, 49.1\%), recursive at mid-high threshold ($\tau=0.50$, 50.9\%). This reveals fundamental difference in optimization dynamics.

\textbf{Threshold crossover}: At $\tau=0.40$, n-way leads (+1.8 pts). At $\tau=0.50$, recursive leads (+4.5 pts). Methods exhibit different parameter landscapes requiring independent tuning.

\textbf{Mechanistic interpretation}: N-way benefits from broad minority definition ($\tau=0.40$) because simultaneous $k$-way optimization can selectively concentrate high-minority tracts while ignoring moderate-minority tracts. Recursive benefits from sharp minority definition ($\tau=0.50$) because hierarchical splits need clear boundaries—ambiguous tracts at $\tau=0.40$ create suboptimal early splits that cascade through tree.

\begin{figure*}[t]
\centering
\includegraphics[width=\textwidth]{figures/fig2_parameter_sensitivity.png}
\caption{Parameter sensitivity analysis showing success rates across weight factors ($\alpha$) and minority thresholds ($\tau$) for both methods. N-way (left) achieves optimal performance at $\alpha=50, \tau=0.40$ (52.3\%), while recursive (right) peaks at $\alpha=100, \tau=0.50$ (56.8\%). Blue boxes indicate best configurations. Color scale: green (higher success) to red (lower success).}
\label{fig:parameter_sensitivity}
\end{figure*}

\subsubsection{Best Parameter Configurations}

Table~\ref{tab:best_configs} presents top-performing configurations for each method.

\begin{table}[h]
\centering
\begin{tabular}{clcc}
\toprule
\textbf{Method} & \textbf{Configuration} & \textbf{States} & \textbf{Success Rate} \\
\midrule
\multirow{3}{*}{N-way}
& $\alpha=50$, $\tau=0.40$ & 23/44 & 52.3\% \\
& $\alpha=100$, $\tau=0.45$ & 23/44 & 52.3\% \\
& $\alpha=5$, $\tau=0.40$ & 23/44 & 52.3\% \\
\midrule
\multirow{3}{*}{Recursive}
& \textbf{$\alpha=100$, $\tau=0.50$} & \textbf{25/44} & \textbf{56.8\%} \\
& $\alpha=50$, $\tau=0.50$ & 24/44 & 54.5\% \\
& $\alpha=10$, $\tau=0.50$ & 23/44 & 52.3\% \\
\bottomrule
\end{tabular}
\caption{Top three parameter configurations by success rate for each method. Recursive's best configuration achieves 56.8\% success rate, exceeding n-way's best (52.3\%) by 4.5 percentage points.}
\label{tab:best_configs}
\end{table}

\textbf{Performance ceiling difference}: Recursive's best configuration (56.8\%) exceeds n-way's best (52.3\%) by 4.5 points, demonstrating \textbf{higher optimization ceiling} with careful parameter tuning. This trades off against n-way's speed advantage.

\textbf{Plateau effect in n-way}: Three distinct n-way configurations tie at 52.3\%, suggesting performance plateau—further parameter exploration unlikely to improve beyond this ceiling.

\textbf{Configuration consensus}: Recursive's top three configurations all use $\tau=0.50$, confirming threshold preference identified in Table~\ref{tab:threshold_sensitivity}. N-way's top configurations cluster around $\tau=0.40$ and $\tau=0.45$.

\begin{figure*}[t]
\centering
\includegraphics[width=\textwidth]{figures/fig5_best_configs.png}
\caption{Top 5 parameter configurations ranked by success rate for each method. N-way (left) achieves 52.3\% peak with $\alpha=50, \tau=0.40$ (gold border). Recursive (right) reaches 56.8\% with $\alpha=100, \tau=0.50$ (gold border). The 4.5 percentage point gap represents recursive's higher optimization ceiling, though requiring more careful parameter tuning.}
\label{fig:best_configs}
\end{figure*}

\subsection{State-by-State Comparison}

\subsubsection{Overall State-Level Performance}

Across 44 multi-district states:
\begin{itemize}
\item \textbf{N-way wins}: 5 states (11.4\%)
\item \textbf{Recursive wins}: 8 states (18.2\%)
\item \textbf{Ties}: 31 states (70.4\%)
\end{itemize}

The high tie rate (70.4\%) supports overall equivalence finding—most states show no method preference.

\subsubsection{Largest Method Advantages}

Table~\ref{tab:state_advantages} shows states with largest performance differences.

\begin{table}[h]
\centering
\begin{tabular}{lccc}
\toprule
\multicolumn{4}{c}{\textbf{States Where N-way Excels}} \\
\textbf{State} & \textbf{N-way} & \textbf{Recursive} & \textbf{Margin} \\
\midrule
Virginia       & 70.0\% & 35.0\% & +35.0 pts \\
Missouri       & 100.0\% & 80.0\% & +20.0 pts \\
Wisconsin      & 65.0\% & 45.0\% & +20.0 pts \\
Maryland       & 100.0\% & 95.0\% & +5.0 pts \\
\midrule
\multicolumn{4}{c}{\textbf{States Where Recursive Excels}} \\
\textbf{State} & \textbf{N-way} & \textbf{Recursive} & \textbf{Margin} \\
\midrule
Connecticut    & 0.0\% & 45.0\% & +45.0 pts \\
Louisiana      & 30.0\% & 55.0\% & +25.0 pts \\
Georgia        & 65.0\% & 80.0\% & +15.0 pts \\
Florida        & 5.0\% & 15.0\% & +10.0 pts \\
Arizona        & 30.0\% & 40.0\% & +10.0 pts \\
\bottomrule
\end{tabular}
\caption{States with largest performance differences between methods. Success rate = average across 20 configurations per method. Shows method-specific advantages exist but concentrate in minority of states.}
\label{tab:state_advantages}
\end{table}

\textbf{Virginia: N-way's strongest advantage} (+35 points) — Virginia has 11 districts, 41.4\% minority, moderately clustered (Washington DC suburbs, Hampton Roads, Richmond). N-way's global $k$-way optimization effectively balances minority concentration across 11 districts. Recursive's early binary splits fragment minorities between northern and southern regions.

\textbf{Connecticut: Recursive's strongest advantage} (+45 points) — Connecticut has 5 districts, 36.8\% minority, dispersed across small cities (Hartford, New Haven, Bridgeport, Stamford). Recursive's hierarchical splits isolate city clusters. N-way's simultaneous 5-way partition struggles to concentrate dispersed minorities without violating compactness.

\textbf{Pattern hypothesis}: N-way excels in moderate-district-count states (8-12 districts) with concentrated minorities. Recursive excels in low-district-count states (4-7 districts) with dispersed minorities. Testing this hypothesis requires additional analysis beyond this paper's scope.

\begin{figure*}[t]
\centering
\includegraphics[width=\textwidth]{figures/fig3_state_comparison.png}
\caption{State-specific method advantages showing top 10 states favoring each method. Positive values indicate n-way advantage (blue); negative values indicate recursive advantage (orange). Most states (31/44) show ties within $\pm$5 percentage points. Notable patterns: n-way excels in states with 8-12 districts and concentrated urban minorities, while recursive performs better in states with 4-7 districts and dispersed minorities.}
\label{fig:state_comparison}
\end{figure*}

\subsubsection{Success Tier Distribution}

Table~\ref{tab:success_tiers} categorizes states by success rate tier for each method.

\begin{table}[h]
\centering
\begin{tabular}{lcc}
\toprule
\textbf{Success Tier} & \textbf{N-way} & \textbf{Recursive} \\
\midrule
100\% (always achieves target) & 17 states & 15 states \\
50-99\% (usually achieves target) & 4 states & 5 states \\
1-49\% (rarely achieves target) & 7 states & 9 states \\
0\% (never achieves target) & 16 states & 15 states \\
\bottomrule
\end{tabular}
\caption{State distribution by success rate tier (average across 20 configurations). Tier counts similar between methods, confirming overall equivalence.}
\label{tab:success_tiers}
\end{table}

\textbf{Bimodal distribution}: Both methods show pronounced bimodality—states either succeed consistently (100\%) or fail consistently (0\%), with few intermediate cases. This indicates VRA compliance depends primarily on state characteristics (minority percentage, geographic clustering, district count) rather than method choice or parameter tuning.

\textbf{Tier stability}: Tier distributions nearly identical between methods (±2 states per tier), confirming statistical equivalence. Methods rarely shift states across success tiers.

\subsection{Runtime Analysis}

\subsubsection{Runtime Distribution}

\begin{table}[h]
\centering
\begin{tabular}{lccccc}
\toprule
\textbf{Method} & \textbf{Mean} & \textbf{Median} & \textbf{Std Dev} & \textbf{Min} & \textbf{Max} \\
\midrule
N-way     & 3.68s & 3.48s & 2.05s & 0.38s & 10.28s \\
Recursive & 11.33s & 10.20s & 7.54s & 0.38s & 39.10s \\
\midrule
Speedup   & 3.08$\times$ & 2.93$\times$ & -- & 1.0$\times$ & 3.80$\times$ \\
\bottomrule
\end{tabular}
\caption{Runtime statistics across 880 configurations per method. Speedup = recursive runtime / n-way runtime. N-way consistently faster across all percentiles.}
\label{tab:runtime_distribution}
\end{table}

\textbf{Consistent speed advantage}: N-way faster at every percentile (minimum, median, mean, maximum). Speedup ranges from 1.0$\times$ (small states, both fast) to 3.80$\times$ (large states like California, Texas).

\textbf{Variance difference}: Recursive shows higher runtime variance (std dev 7.54s vs 2.05s), indicating less predictable performance. Large states with complex geometries show extreme slowdown under recursive (39.10s maximum).

\textbf{Scaling behavior}: Runtime scales with state complexity (district count, tract count, geographic irregularity). N-way's single $k$-way optimization scales polynomially. Recursive's $\log_2(k)$ split depth partially offsets this, but each split operates on full remaining graph, leading to worse constant factors.

\begin{figure*}[t]
\centering
\includegraphics[width=\textwidth]{figures/fig4_runtime_comparison.png}
\caption{Runtime distribution comparison. Left: box plots showing median, quartiles, and outliers. N-way (blue) has tighter distribution and lower median (3.68s) compared to recursive (orange, 11.33s). Right: histograms showing frequency distributions. N-way exhibits 67.5\% speed advantage, crucial for ensemble generation and interactive applications.}
\label{fig:runtime_distribution}
\end{figure*}

\subsubsection{Runtime vs Success Rate Trade-off}

\textbf{Efficiency analysis}: To gain 0.8 percentage points (5 additional state successes out of 880 trials), recursive requires 208\% more computation time. Cost per marginal success: 1.53s. For ensemble generation (10,000 maps), this translates to 15,300s (4.25 hours) additional computation for ~8 additional state successes.

\textbf{Practical implication}: For applications prioritizing speed (ensemble generation, rapid prototyping, real-time optimization), n-way dominates. For applications prioritizing maximum success rate (single carefully-tuned plan, official redistricting), recursive's 0.8 point advantage may justify runtime cost—but only if parameter tuning can reliably access the 56.8\% ceiling.

\subsection{Summary of Key Results}

\textbf{Result 1: Statistical Equivalence} — Both methods achieve ~48\% success rate with negligible difference (0.8 points, $p=0.634$, $d=-0.018$). Architectural choice does not substantially affect VRA compliance.

\textbf{Result 2: Speed-Quality Trade-off} — N-way is 67.5\% faster (3.68s vs 11.33s). Recursive achieves slightly higher average (48.3\% vs 47.5\%) and substantially higher ceiling (56.8\% vs 52.3\% best config). Trade-off depends on application requirements.

\textbf{Result 3: Different Parameter Landscapes} — N-way optimizes at $\tau=0.40$, recursive at $\tau=0.50$. Weight preferences also differ, requiring independent parameter tuning per method.

\textbf{Result 4: State-Specific Advantages} — Methods tie in 70.4\% of states. N-way excels in Virginia-type states (moderate districts, concentrated minorities). Recursive excels in Connecticut-type states (few districts, dispersed minorities).

\textbf{Result 5: Bimodal Success Distribution} — Both methods show strong bimodality (100\% or 0\% success). State characteristics (demographics, geography) dominate over method choice, explaining overall equivalence despite architectural differences.
